\SourceLeaf{230}%
\Article{5}
\indent Duae solutiones in art. praec. inventae manifesto ad puncta superficiei sphaericae opposita, sive ad directiones oppositas referuntur, quod cum rei natura quadrat, quum normalem ad utramvis plagam superficiei curvae ducere liceat. Quodsi duas plagas, superficiei contiguas, inter se distinguere, alteramque exteriorem alteram interiorem vocare placet, etiam utrique normali suam solutionem rite tribuere licebit adiumento theorematis in art. 2 (VII) evoluti, simulatque criterium stabilitum est ad plagam alteram ab altera distinguendam.

In methodo prima tale criterium petendum erit a signo valoris quantitatis \(W\). Scilicet generaliter loquendo superficies curva eas spatii partes, in quibus \(W\) valorem positivum obtinet, ab iis dirimet, in quibus valor ipsius \(W\) fit negativus. E theoremate illo vero facile colligitur, si \(W\) valorem positivum obtineat versus plagam exteriorem, normalisque extrorsum ducta concipiatur, solutionem priorem adoptandam esse. Ceterum in quovis casu facile diiudicabitur, utrum per superficiem integram eadem regula respectu signi ipsius \(W\) valeat, an pro diversis partibus diversae: quamdiu coëfficientes \(P,Q,R\) valores finitos habent, nec simul omnes tres evanescunt, lex continuitatis vicissitudinem vetabit.

Si methodum secundam sequimur, in superficie curva duo systemata linearum curvarum concipere possumus, alterum, pro quo \(p\) est variabilis, \(q\) constans; alterum, pro quo \(q\) variabilis, \(p\) constans: situs mutuus harum linearum respectu plagae exterioris decidere debet, utram solutionem adoptare oporteat. Scilicet quoties tres lineae, puta ramus lineae prioris systematis a puncto \(A\) proficiscens crescente \(p\), ramus posterioris systematis a puncto \(A\) egrediens crescente \(q\), atque normalis versus plagam exteriorem ducta \textit{similiter} iacent, ut, inde ab origine abscissarum, axes ipsarum \(x,y,z\) resp. (e. g. si tum e tribus lineis illis, tum e tribus his, prima sinistrorsum, secunda dextrorsum, tertia sursum directa concipi potest), solutio prima adoptari debet; quoties autem situs mutuus trium linearum oppositus est situi mutuo axium ipsarum \(x,y,z\), solutio secunda valebit.

In methodo tertia dispiciendum est, utrum, dum \(z\) incrementum positivum accipit, manentibus \(x\) et \(y\) invariatis, transitus fiat versus plagam exteriorem an interiorem. In casu priore, pro normali extrorsum directa, solutio prima valet, in posteriore secunda.
